\documentclass[12pt]{letter}
\usepackage{geometry}
\geometry{margin=1in}
\usepackage{hyperref}

\begin{document}

\begin{letter}{%
The Editors\\
\emph{Theory and Decision}\\
Springer}

\opening{Dear Editors,}

I am pleased to submit the manuscript ``Geometric Economics: Heuristic Bounded Rationality and the Bond Geodesic Equilibrium'' for consideration as an Original Article in \emph{Theory and Decision}.

The paper develops a mathematical framework in which economic decision-making is modeled as pathfinding on a multi-dimensional decision manifold rather than scalar utility maximization.  Agents evaluate alternatives on a nine-dimensional weighted simplicial complex---the economic decision complex---whose dimensions include monetary cost, moral constraint, social norm, fairness, autonomy, trust, identity, institutional legitimacy, and epistemic status.  The agent's decision is modeled as $A^*$ search with evaluation function $f(n) = g(n) + h(n)$, where $g(n)$ is accumulated economic cost (Kahneman's System~2) and $h(n)$ is the moral-heuristic estimate (System~1).

The paper makes four principal contributions:
\begin{enumerate}
    \item \textbf{Heuristic Bounded Rationality:} A proof that moral and cultural heuristics are admissible (or $\varepsilon$-admissible) $A^*$ search heuristics, not cognitive biases; together with a formal treatment of the admissibility spectrum from strictly optimal to bounded-suboptimal search.
    \item \textbf{The Scalar Irrecoverability Theorem:} A proof (via Brouwer's invariance of dimension) that no continuous function can losslessly project the nine-dimensional decision space to a scalar utility, establishing that information loss is inherent in classical models.
    \item \textbf{The Bond Geodesic Equilibrium (BGE):} A generalization of Nash equilibrium for agents who optimize on the full manifold.  BGE is shown to nest Nash equilibrium as a special case via an augmented game construction, with existence guaranteed by Nash's theorem on the augmented game, and uniqueness conditions given by an explicit contraction lemma.
    \item \textbf{Geometric Behavioral Economics:} Systematic derivation of prospect-theoretic phenomena (loss aversion, framing effects, reference dependence, hyperbolic discounting) as geometric properties of the decision manifold's metric.
\end{enumerate}

This work is relevant to \emph{Theory and Decision} for three reasons.  First, it provides a new mathematical foundation for bounded rationality that formalizes Simon's program using $A^*$ pathfinding, connecting decision theory to computational search theory.  Second, it offers a unified derivation of behavioral economics anomalies as theorems from manifold structure, rather than cataloguing them as separate phenomena.  Third, the BGE--Nash nesting result and the contraction conditions for uniqueness contribute directly to game-theoretic foundations.

The nine-dimensional moral manifold that underpins the economic decision complex has been empirically validated across 20,030 texts and 11 languages (Bond, 2026), with independent replication across 6 additional languages (Thiele, 2026).  The paper proposes six falsifiable predictions that distinguish the framework from both classical and behavioral economics.

The manuscript has not been published previously and is not under consideration elsewhere.

Thank you for considering this submission.

\closing{Sincerely,\\[2em]
Andrew H. Bond\\
Senior Lecturer, Department of Computer Engineering\\
San Jos\'{e} State University\\
\href{mailto:andrew.bond@sjsu.edu}{andrew.bond@sjsu.edu}\\
ORCID: 0009-0003-2599-6158}

\end{letter}
\end{document}
